%---------------------------------------------------------------------
% Part one: the short-definition question.
%
% This is set out as a glossary rather than as a numbered question,
% because in two of the four papers it IS question one and carries a
% full 20 marks on its own -- ten one-line answers. Every term that was
% actually asked in 2013 or 2021 is here, with the paper marked; the
% rest are the terms that sit alongside them in the same chapter and
% are the obvious next candidates.
%---------------------------------------------------------------------

\section*{حصہ اول --- مختصر اصطلاحات}
\markboth{مختصر اصطلاحات}{}

\noindent
بہار \textenglish{2013} اور خزاں \textenglish{2021} دونوں میں \textbf{پہلا
سوال یہی تھا اور لازمی تھا}: دس اصطلاحات، ہر ایک کا مختصر جواب، کل بیس نمبر۔
یہ پرچے کا سب سے سستا سوال ہے --- ہر جواب صرف دو تین سطروں کا ہے۔ ذیل میں وہ
تمام اصطلاحات ہیں جو واقعی پوچھی گئیں (پرچے کا نام ساتھ درج ہے) اور وہ جو انہی
ابواب سے اگلی بار متوقع ہیں۔

\subsection*{قرآن و حدیث کی اصطلاحات}

\begin{center}
\begin{tabular}{@{}p{3.3cm} p{11.1cm}@{}}
\toprule
\textbf{محکمات و متشابہات}
  & \textbf{محکمات}: واضح المعنی آیات جن میں تاویل کی گنجائش نہیں --- یہی
    ''اُمُّ الْکِتَاب'' یعنی کتاب کی اصل ہیں۔ \textbf{متشابہات}: جن کا حقیقی
    مفہوم اللہ کے سوا کوئی نہیں جانتا۔ بنیاد: \ar{مِنْهُ آيَاتٌ مُّحْكَمَاتٌ
    هُنَّ أُمُّ الْكِتَابِ وَأُخَرُ مُتَشَابِهَاتٌ} (آل عمران، \textenglish{7})۔
    {\small\color{bmgrey}[خزاں \textenglish{2021}]} \\
\textbf{حدیث قدسی}
  & وہ حدیث جس کا \emph{مضمون} اللہ کی طرف سے ہو اور \emph{الفاظ} نبی کریم
    ﷺ کے، اور جو ''\ar{قَالَ اللهُ تَعَالَىٰ}'' کے ساتھ روایت ہو۔ قرآن سے فرق:
    قرآن لفظاً و معناً منزل ہے، متواتر ہے، نماز میں اس کی تلاوت ہوتی ہے اور وہ
    معجزہ ہے --- حدیث قدسی ان میں سے کسی وصف میں شریک نہیں۔
    {\small\color{bmgrey}[بہار \textenglish{2013}]} \\
\textbf{صحاح ستہ}
  & حدیث کی چھ مستند کتابیں: صحیح بخاری (م \textenglish{256}ھ)، صحیح مسلم
    (\textenglish{261}ھ)، سنن ابن ماجہ (\textenglish{273}ھ)، سنن ابو داؤد
    (\textenglish{275}ھ)، جامع ترمذی (\textenglish{279}ھ)، سنن نسائی
    (\textenglish{303}ھ)۔ پہلی دو ''صحیحین'' کہلاتی ہیں۔
    {\small\color{bmgrey}[بہار \textenglish{2013}]} \\
\textbf{مباہلہ}
  & فریقین کا اپنے اہل و عیال سمیت جمع ہو کر یہ دعا کرنا کہ جھوٹے پر اللہ کی
    لعنت ہو۔ آل عمران، \textenglish{61} میں نجران کے عیسائی وفد کو یہی دعوت دی
    گئی؛ وہ سامنے آنے سے انکار کر کے جزیہ پر راضی ہو گئے۔ \\
\bottomrule
\end{tabular}
\end{center}

\subsection*{عقائد و اخلاق}

\begin{center}
\begin{tabular}{@{}p{3.3cm} p{11.1cm}@{}}
\toprule
\textbf{تقدیر مبرم} \newline \textbf{اور معلق}
  & \textbf{مبرم}: اللہ کا اٹل فیصلہ جو بدلتا نہیں۔ \textbf{معلق}: جو کسی سبب
    --- دعا، صدقہ، صلہ رحمی --- سے مشروط ہے اور بدل سکتا ہے۔ بنیاد:
    \ar{يَمْحُو اللهُ مَا يَشَاءُ وَيُثْبِتُ وَعِندَهُ أُمُّ الْكِتَابِ}
    (الرعد، \textenglish{39})، اور ''\ar{لَا يَرُدُّ الْقَضَاءَ إِلَّا
    الدُّعَاءُ}'' (ترمذی)۔
    {\small\color{bmgrey}[بہار \textenglish{2013}، خزاں \textenglish{2021} ---
    دونوں میں]} \\
\textbf{کبر}
  & نبی کریم ﷺ نے خود اس کی تعریف فرمائی: ''\ar{الْكِبْرُ بَطَرُ الْحَقِّ
    وَغَمْطُ النَّاسِ}'' --- حق کو ٹھکرا دینا اور لوگوں کو حقیر جاننا (مسلم)۔
    اسی حدیث میں ہے کہ جس کے دل میں ذرہ برابر کبر ہو وہ جنت میں داخل نہ ہو گا۔
    {\small\color{bmgrey}[بہار \textenglish{2013}]} \\
\textbf{حسن خلق}
  & لوگوں کے ساتھ خوش اخلاقی، نرمی، ایفائے عہد اور ایذا سے بچنا۔ ''\ar{الْبِرُّ
    حُسْنُ الْخُلُقِ}'' --- نیکی خوش خلقی ہی کا نام ہے (مسلم)، اور
    ''\ar{أَكْمَلُ الْمُؤْمِنِينَ إِيمَانًا أَحْسَنُهُمْ خُلُقًا}'' (ترمذی)۔
    {\small\color{bmgrey}[بہار \textenglish{2013}]} \\
\textbf{الحب فی اللہ}
  & کسی سے محض اللہ کی رضا کے لیے محبت کرنا --- نہ نسب کی بنا پر، نہ مال یا
    مفاد کی۔ جن سات آدمیوں کو قیامت کے دن عرش کا سایہ ملے گا ان میں
    ''\ar{وَرَجُلَانِ تَحَابَّا فِي اللهِ اجْتَمَعَا عَلَيْهِ وَتَفَرَّقَا
    عَلَيْهِ}'' بھی ہیں (بخاری و مسلم)۔
    {\small\color{bmgrey}[بہار \textenglish{2013}، خزاں \textenglish{2021}]} \\
\textbf{رحمت و شفقت} \newline \textbf{کی حد}
  & رحمت عام ہے، مگر اس کی حد وہاں ہے جہاں \textbf{اللہ کی حد} شروع ہوتی ہے۔
    قرآن: \ar{وَلَا تَأْخُذْكُم بِهِمَا رَأْفَةٌ فِي دِينِ اللهِ} (النور،
    \textenglish{2})۔ جب حضرت اسامہؓ نے چوری کرنے والی خاتون کی سفارش کی تو
    فرمایا: ''\ar{أَتَشْفَعُ فِي حَدٍّ مِنْ حُدُودِ اللهِ؟}'' (بخاری و مسلم)۔
    یعنی نرمی جرم کی سزا معاف کرانے کا ذریعہ نہیں بن سکتی۔
    {\small\color{bmgrey}[خزاں \textenglish{2021}]} \\
\textbf{اتباع اور} \newline \textbf{اطاعت کا فرق}
  & \textbf{اطاعت} حکم کی تعمیل ہے --- جو کہا گیا وہ کر دیا۔ \textbf{اتباع}
    قدم بہ قدم پیروی ہے --- قول، فعل اور طرزِ زندگی سب میں۔ اتباع اطاعت سے
    وسیع تر ہے۔ اسی لیے محبت کی دلیل اتباع رکھی گئی: \ar{قُلْ إِن كُنتُمْ
    تُحِبُّونَ اللهَ فَاتَّبِعُونِي} (آل عمران، \textenglish{31})۔
    {\small\color{bmgrey}[خزاں \textenglish{2021}]} \\
\textbf{تشمیت}
  & تشمیتِ عاطس: چھینکنے والا جب ''\ar{الْحَمْدُ لِلهِ}'' کہے تو سننے والے کا
    ''\ar{يَرْحَمُكَ اللهُ}'' کہنا، اور پھر چھینکنے والے کا ''\ar{يَهْدِيكُمُ
    اللهُ وَيُصْلِحُ بَالَكُمْ}'' کہنا (بخاری)۔ یہ مسلمان کے چھ حقوق میں سے ایک
    ہے۔ لفظی مفہوم: دعائے خیر دینا۔
    {\small\color{bmgrey}[بہار \textenglish{2013}]} \\
\textbf{دینیات}
  & وہ علم جو دینِ اسلام کے بنیادی مآخذ (قرآن و سنت) اور اس کے شعبوں ---
    عقائد، عبادات، اخلاق، معاملات اور سیرت و تاریخ --- کا منظم مطالعہ کرتا ہے۔
    یہی اس کورس کا موضوع ہے۔
    {\small\color{bmgrey}[بہار \textenglish{2013}]} \\
\bottomrule
\end{tabular}
\end{center}

\clearpage

\subsection*{عبادات کی اصطلاحات}

\begin{center}
\begin{tabular}{@{}p{3.3cm} p{11.1cm}@{}}
\toprule
\textbf{تکبیراتِ تشریق}
  & نو ذی الحجہ کی فجر سے تیرہ ذی الحجہ کی عصر تک ہر فرض نماز کے بعد ایک بار
    بلند آواز سے کہنا: ''\ar{اللهُ أَكْبَرُ اللهُ أَكْبَرُ، لَا إِلٰهَ إِلَّا
    اللهُ وَاللهُ أَكْبَرُ، اللهُ أَكْبَرُ وَلِلهِ الْحَمْدُ}''۔ کل
    \textbf{تیئس} نمازیں بنتی ہیں۔ حنفیہ کے نزدیک واجب۔ بنیاد: \ar{وَاذْكُرُوا
    اللهَ فِي أَيَّامٍ مَّعْدُودَاتٍ} (البقرہ، \textenglish{203})۔
    {\small\color{bmgrey}[خزاں \textenglish{2021}]} \\
\textbf{کسوف و خسوف}
  & \textbf{کسوف} سورج گرہن، \textbf{خسوف} چاند گرہن۔ سورج گرہن کی دو رکعت
    نماز باجماعت سنت ہے --- بغیر اذان و اقامت کے۔ چاند گرہن کی نماز حنفیہ کے
    نزدیک انفرادی پڑھی جائے۔ حدیث: ''\ar{إِنَّ الشَّمْسَ وَالْقَمَرَ آيَتَانِ
    مِنْ آيَاتِ اللهِ، لَا يَنْخَسِفَانِ لِمَوْتِ أَحَدٍ وَلَا لِحَيَاتِهِ}''
    (بخاری و مسلم) --- یہ ارشاد حضرت ابراہیمؓ کی وفات کے دن ہوا، جب لوگوں نے
    گرہن کو اس سے جوڑ دیا تھا۔
    {\small\color{bmgrey}[خزاں \textenglish{2021}]} \\
\textbf{شفق}
  & غروبِ آفتاب کے بعد مغربی افق پر باقی رہنے والی روشنی۔ اس کے غائب ہوتے ہی
    مغرب کا وقت ختم اور عشاء کا شروع ہوتا ہے۔ امام ابو حنیفہؒ کے نزدیک یہ
    \emph{سفیدی} ہے، صاحبین اور جمہور کے نزدیک \emph{سرخی} --- اسی اختلاف سے
    عشاء کا وقت آگے پیچھے ہوتا ہے۔ \\
\textbf{سجدہ سہو}
  & بھول کر کوئی واجب چھوٹ جانے یا اس میں تاخیر ہونے پر آخری قعدے میں تشہد کے
    بعد ایک طرف سلام پھیر کر دو سجدے کرنا، پھر تشہد، درود اور دعا پڑھ کر سلام
    پھیرنا۔ یاد رہے: \textbf{فرض} چھوٹ جائے تو سجدہ سہو کافی نہیں --- نماز
    دوبارہ پڑھنی پڑے گی۔ \\
\bottomrule
\end{tabular}
\end{center}

\subsection*{تاریخ اور مالیات کی اصطلاحات}

\begin{center}
\begin{tabular}{@{}p{3.3cm} p{11.1cm}@{}}
\toprule
\textbf{خراج}
  & غیر مسلم رعایا کی \textbf{زرعی زمین} پر مقرر کیا جانے والا لگان۔ دو قسمیں:
    \textbf{خراجِ موظف} --- زمین پر سالانہ مقررہ رقم، خواہ پیداوار کم ہو یا
    زیادہ؛ اور \textbf{خراجِ مقاسمہ} --- پیداوار کا مقررہ حصہ (مثلاً چوتھائی
    یا پانچواں)۔ حضرت عمرؓ نے عراق و شام کی مفتوحہ زمینوں پر اس کا باقاعدہ
    نظام قائم کیا۔
    {\small\color{bmgrey}[بہار \textenglish{2013}، خزاں \textenglish{2021} ---
    دونوں میں]} \\
\textbf{جزیہ}
  & ذمی غیر مسلموں کے \emph{بالغ، آزاد، صاحبِ حیثیت مردوں} پر جان و مال کے
    تحفظ کے بدلے سالانہ ٹیکس۔ عورتوں، بچوں، بوڑھوں، معذوروں، راہبوں اور
    فقیروں پر نہیں۔ یہ فوجی خدمت سے استثنا کا معاوضہ ہے۔ بنیاد: \ar{حَتَّىٰ
    يُعْطُوا الْجِزْيَةَ عَن يَدٍ} (التوبہ، \textenglish{29})۔ \\
\textbf{عشر}
  & مسلمان کی زرعی پیداوار پر زکوٰۃ۔ بارانی یا چشمے سے سیراب زمین پر
    \textbf{دسواں حصہ}، مصنوعی آبپاشی والی زمین پر \textbf{بیسواں حصہ}۔ حدیث:
    ''\ar{فِيمَا سَقَتِ السَّمَاءُ وَالْعُيُونُ الْعُشْرُ، وَمَا سُقِيَ
    بِالنَّضْحِ نِصْفُ الْعُشْرِ}'' (بخاری)۔ \\
\textbf{عَنوہ اور صلحاً}
  & \textbf{عَنوۃً} فتح ہونے والا علاقہ وہ ہے جو \emph{جنگ اور طاقت} سے حاصل
    ہو؛ \textbf{صلحاً} وہ جو \emph{معاہدے} سے۔ فرق محض تاریخی نہیں --- زمین کی
    ملکیت، خراج کی نوعیت اور اہلِ علاقہ کے حقوق سب اسی بنیاد پر طے ہوتے ہیں۔
    {\small\color{bmgrey}[خزاں \textenglish{2021}]} \\
\textbf{سفاح ثانی}
  & عباسی خلیفہ \textbf{المعتضد باللہ} (\textenglish{279--289}ھ) کا لقب۔ پہلے
    ''سفاح'' ابو العباس عبد اللہ تھے، یعنی خلافتِ عباسیہ کے بانی و پہلے خلیفہ۔
    المعتضد کو یہ لقب اس لیے ملا کہ اس نے زوال کے دہانے پر پہنچی ہوئی عباسی
    طاقت کو سختی سے دوبارہ مستحکم کیا۔
    {\small\color{bmgrey}[بہار \textenglish{2013}]} \\
\textbf{تاتاری}
  & وسطی ایشیا کے منگول قبائل، جنہیں چنگیز خان (م \textenglish{624}ھ) نے متحد
    کیا۔ اسی کے پوتے ہلاکو خان نے \textenglish{656}ھ میں بغداد تباہ کر کے
    خلافتِ عباسیہ ختم کی۔ دلچسپ بات یہ کہ ایک صدی کے اندر بہت سے منگول خود
    مسلمان ہو گئے۔
    {\small\color{bmgrey}[بہار \textenglish{2013}]} \\
\textbf{وادی حنین}
  & مکہ اور طائف کے درمیان واقع ایک وادی۔ یہاں \textenglish{8}ھ میں، فتحِ مکہ
    کے فوراً بعد، ہوازن اور ثقیف کے قبائل سے غزوہ حنین ہوا۔ ابتدا میں مسلمانوں
    کے قدم اکھڑ گئے تھے: \ar{وَيَوْمَ حُنَيْنٍ إِذْ أَعْجَبَتْكُمْ
    كَثْرَتُكُمْ فَلَمْ تُغْنِ عَنكُمْ شَيْئًا} (التوبہ، \textenglish{25})۔
    {\small\color{bmgrey}[خزاں \textenglish{2021}]} \\
\bottomrule
\end{tabular}
\end{center}

\medskip
\begin{nukta}[اس سوال میں نمبر کیسے پورے ہوں]
دس اصطلاحات، بیس نمبر --- یعنی \textbf{فی اصطلاح دو نمبر}، اور وقت صرف بیس
پچیس منٹ۔ لمبا مضمون لکھنے کی ضرورت نہیں۔ ہر جواب میں \textbf{ایک سطر تعریف}
اور \textbf{ایک حوالہ} (آیت، حدیث یا سنہ) کافی ہے۔ حوالہ ہی وہ چیز ہے جو دو
میں سے دوسرا نمبر دلاتی ہے۔
\end{nukta}
